\documentclass[12pt]{article}
\usepackage[utf8]{inputenc}
\usepackage[T1]{fontenc}
\usepackage{hyperref}
\title{Sudanese Arabic Language Resources for K-12 Curriculum Development}
\author{Compiled by Claude Sonnet 5.5 (claude-sonnet-5-5) for Andr\'as Kornai}
\date{}
\begin{document}
\maketitle

\section{Introduction}
Sudanese Arabic (ISO 639-3: apd) is the colloquial Arabic of Sudan and South Sudan. Search turned up only learner-oriented materials aimed at foreigners; the dictionary's own description notes that the spoken language is largely unwritten. Schooling in Sudan uses Modern Standard Arabic, and no Sudanese Arabic curriculum materials for schools were found.

\section{Primary Educational Resources}
No school textbooks or primers in Sudanese colloquial Arabic were found. The learner coursebook below is a language-learning text rather than a school resource.
\begin{itemize}
\item Sudanese Colloquial Arabic for Beginners (Janet Persson, Ahmad Hussein, Andrew M. Persson; SIL, 2nd ed.): dialogues, drills, grammar notes, and per-lesson vocabulary. Vendor listing: \url{https://kenyaclc.org/shop/sudanese-colloquial-arabic-for-beginners/}
\end{itemize}

\section{Linguistic and Reference Materials}
\begin{itemize}
\item Sudanese Arabic--English / English--Sudanese Arabic: A Concise Dictionary (Rianne Tamis and Janet Persson): \url{https://kenyaclc.org/shop/sudanese-arabic-english-english-sudanese-arabic-a-concise-dictionary/}
\item Glottolog bibliographic entry for the Persson, Hussein and Persson 1980 coursebook: \url{https://glottolog.org/resource/reference/id/133773}
\end{itemize}

\section{Religious and Community Resources}
None found in search results.

\section{Digital and Audio Resources}
No audio, video or corpus resources were found in search results.

\section{Curriculum Development Context}
Only two learner-oriented reference works were found, both for outsiders learning the dialect. Native-language curriculum work would need Sudanese teachers and linguists, and a decision on how to relate a mostly unwritten colloquial to the standard Arabic used in schools. Little else could be verified here.

\end{document}